% -----------------------------------------------------------------------------
\section{Background, Related Work, and Research Inputs}
\label{sec:background}
% -----------------------------------------------------------------------------

MoonAI sits at the intersection of neuroevolution, artificial life simulation, GPU computing,
and interactive experiment tooling. The final design was shaped by both academic background
and practical engineering references.

\subsection{Neuroevolution and Predator-Prey Simulation}

The underlying learning approach is NeuroEvolution of Augmenting Topologies (NEAT), which
evolves both neural network weights and network structure over time \cite{neat2002}. NEAT is
relevant to MoonAI because the project is concerned with adaptation under pressure, not only
with parameter fitting inside a fixed model architecture.

Predator-prey worlds are a natural fit for this kind of study. They create coevolutionary
pressure, local competition, and non-stationary behavior: the quality of a strategy depends on
what other agents are doing at the same time. This makes the environment useful for observing
emergent behavior, species diversity, and topology growth under changing conditions.

\subsection{Why a Synthetic Environment Matters}

The project deliberately prioritizes a synthetic environment over a pre-collected dataset.
That decision has several research advantages:

\begin{itemize}[nosep]
  \item the environment generates its own behavior traces and metrics,
  \item experiments can be rerun with fixed seeds and controlled parameter sweeps,
  \item selective pressure is explicit and configurable,
  \item the benchmark can be extended without collecting external labeled data.
\end{itemize}

This does not make synthetic environments universally superior to real data. Instead, it makes
them appropriate for a different class of questions: questions about adaptation, dynamics,
policy diversity, and algorithmic behavior under repeated interaction.

\subsection{Related Engineering Inputs}

The final implementation was also shaped by practical design inputs rather than by a single
prior system. The team relied on official documentation and mature open-source ecosystems for:

\begin{itemize}[nosep]
  \item Rust systems programming and package management \cite{rust},
  \item CUDA kernel programming and host-device orchestration \cite{cuda},
  \item immediate-mode desktop UI and GPU rendering through \texttt{egui}, \texttt{eframe},
    and \texttt{wgpu} \cite{egui,wgpu},
  \item Lua-based runtime configuration through \texttt{mlua} \cite{mlua},
  \item JSON serialization and persistence through \texttt{serde} \cite{serde},
  \item post-run analysis and report generation through Pandas, Matplotlib, NetworkX, and
    Jinja2 \cite{pandas,matplotlib,networkx,jinja2}.
\end{itemize}

These sources were used to understand component behavior, data handling, rendering patterns,
and reporting workflows. They informed the project architecture, but the final integration and
runtime behavior are specific to MoonAI.

\subsection{Use of Library and Internet Resources}

Library and Internet resources were used in four main ways during the course:

\begin{enumerate}[nosep]
  \item to study background concepts such as NEAT and evolutionary computation,
  \item to evaluate technology choices for GPU execution, UI rendering, serialization, and
    report generation,
  \item to check component APIs, version compatibility, and workflow constraints,
  \item to document and visualize the final system using GitHub Pages, LaTeX, and Mermaid
    diagrams \cite{mermaid}.
\end{enumerate}

The project therefore depends on both academic references and engineering documentation. The
academic literature explains why the problem is interesting; the technical references explain
how the system can be implemented reliably.
